How much of a language model's parameter budget is spent on standard categorical and conditional knowledge --- the kind a competent reader can rattle off about peanuts, Georgia, or symphony orchestras? Capacity scaling laws tell us how many bits a transformer can store in total; mechanistic interpretability tells us where individual facts live; probing benchmarks tell us how many facts a model knows. None of these threads directly answers the share question. We close that gap by combining log-likelihood probing on \bear and \taxi with the \citet{allenzhu2024knowledge} \(2\) bits-per-parameter capacity ceiling. Across \pythia models from \(160\)M to \(2.8\)B parameters, standard categorical and relational knowledge consumes only \(6\times 10^{-5}\%\) to \(2 \times 10^{-4}\%\) of the bit budget --- roughly one part in a million of stored capacity. The fraction shrinks monotonically with model size from \(410\)M onward (slope \(-1.5 \times 10^{-4}\) percentage-points per log-decade, \(p=0.05\)), with bootstrap \(95\%\) confidence intervals at \(1\)B and \(2.8\)B disjoint from those at smaller scales. Standard knowledge is therefore not a binding capacity constraint at modern scales: the vast majority of an LLM's parameters encode something else --- long-tail facts, linguistic competence, reasoning circuitry, and raw memorization.
